\section{Computational Methods}\label{sec:solvers}

\subsection{CVXPY Reference Implementation}

We use CVXPY \cite{diamond2016} as a ground-truth reference: the minimum variance
objective maps directly to a least-squares norm, which CVXPY dispatches to the Clarabel
interior-point solver \cite{clarabel2024} as a second-order cone program.

\begin{algorithm}[ht]
\caption{CVXPY reference solver}\label{alg:cvxpy}
\begin{algorithmic}[1]
\Require Stacked matrix $\tilde{X} \in \mathbb{R}^{(T+n_R)\times n}$;\; return tilt $\rho\hat\mu$;\;
         balance system $B$, $c$
\Ensure Weight vector $w$ with $Bw = c$, $w \geq \mathbf{0}$;\; interior-point iteration count $k$
\State Declare $w \in \mathbb{R}^n$ as a CVXPY variable
\State Minimise $\|\tilde{X}w\|^2 - \rho\hat\mu^\top w$
  \Statex\hspace{1.5em} subject to $Bw = c$,\; $w \geq \mathbf{0}$
\State Dispatch to Clarabel interior-point solver \cite{clarabel2024}
\State \Return $w$,\; $k$
\end{algorithmic}
\end{algorithm}

Interior-point methods require $O(n^3)$ work per iteration for dense problems (as here,
where the covariance matrix is fully dense) plus problem-construction overhead.
On the S\&P~500 problem ($n = 494$) this amounts to roughly 1.9\,s (Table~\ref{tab:sp500}),
making CVXPY unsuitable for production use but ideal as a correctness reference.

All solvers in this paper (CVXPY, CG, and proximal gradient) solve the
same optimization problem; on the S\&P~500 problem the maximum absolute deviation in
portfolio weights between CG and CVXPY (Clarabel) is below $10^{-5}$, confirming
correctness. The comparison is purely computational. CVXPY and Clarabel are general-purpose tools: the same
code handles additional convex constraints (sector limits, turnover bounds, factor-exposure
constraints) without modification. The matrix-free methods accommodate any linear balance
system $Bw = c$ through the Schur elimination of Section~\ref{sec:pd}, but \emph{inequality}
constraints beyond $w \geq 0$ fall outside the reduction and would require rework;
practitioners with such requirements may find a specialized QP solver a more appropriate
baseline.

\subsection{CG Method}\label{ssec:cg}

The inner solver $f$ for Algorithm~\ref{alg:elim} is conjugate gradients applied to the
$n_a \times n_a$ SPD reduction derived in Section~\ref{sec:pd}: $p$ or $p+1$ CG solves
sharing one operator, followed by the $p \times p$ Schur system for the multipliers.
The action of
$\hat\Sigma_{\mathrm{LW},a}$ on a vector requires only two matrix-vector products with
$\tilde{X}_a$; the matrix itself is never formed. Algorithm~\ref{alg:cg} shows the
budget instantiation ($p = 1$) implemented in the companion package and run in every
experiment of Section~\ref{sec:results}.

\begin{algorithm}[ht]
\caption{CG inner step (inner solver $f$ for Algorithm~\ref{alg:elim}), budget case}\label{alg:cg}
\begin{algorithmic}[1]
\Require Active-asset index set $\mathrm{active}$;\; stacked matrix $\tilde{X}$;\;
         return-tilt weight $\rho$;\; expected returns $\hat\mu$
\Ensure Active-asset weights $w_{\mathrm{a}}$;\; CG iteration count $k$
\Function{CgSolve}{active}
\State Let $\tilde{X}_{\mathrm{a}} \leftarrow \tilde{X}_{:,\mathrm{active}}$,\;
       $n_a \leftarrow |\mathrm{active}|$
\State Build matrix-free operator $\mathcal{S}$:\;
    $v \;\mapsto\; \tilde{X}_{\mathrm{a}}^\top(\tilde{X}_{\mathrm{a}}\,v)$
    \hfill\textit{($\hat\Sigma_{\mathrm{LW},a}$ matrix-free, $O(n_a T)$)}
\If{$\rho = 0$}
  \State Solve $\mathcal{S}\,v_1 = \mathbf{1}$ via CG
  \State \Return $v_1\,/\,(\mathbf{1}^\top v_1)$,\; $k$
\Else
  \State Solve $\mathcal{S}\,v_1 = \mathbf{1}$ and $\mathcal{S}\,v_2 = \hat\mu_{\mathrm{active}}$ via CG
  \State $\hat\lambda \leftarrow 2(1 - \tfrac{\rho}{2}\mathbf{1}^\top v_2)\,/\,(\mathbf{1}^\top v_1)$
    \hfill\textit{(budget multiplier, exact)}
  \State \Return $\tfrac{\hat\lambda}{2}v_1 + \tfrac{\rho}{2}v_2$,\; $k_1+k_2$
\EndIf
\EndFunction
\end{algorithmic}
\end{algorithm}

When an asset is dropped in the primal step, $\mathcal{S}$ is rebuilt for the reduced active
set and CG restarts from scratch on the smaller system.

\paragraph{Convergence speed}
\begin{proposition}[CG convergence {\cite{greenbaum1997}}]\label{prop:cg}
  Let $\hat\Sigma_a$ be SPD with spectral condition number
  $\kappa = \kappa(\hat\Sigma_{\mathrm{LW},a})$, and let $v_k$ denote the $k$-th CG
  iterate for $\hat\Sigma_a v = \mathbf{1}$. Then
  \[
    \frac{\lVert e_k\rVert_{\hat\Sigma_a}}{\lVert e_0\rVert_{\hat\Sigma_a}}
    \;\leq\;
    2\!\left(\frac{\sqrt{\kappa}-1}{\sqrt{\kappa}+1}\right)^{\!k},
  \]
  where $\lVert e\rVert_{\hat\Sigma_a} = (e^\top\hat\Sigma_a\,e)^{1/2}$ is the
  $\hat\Sigma_a$-energy norm of the error $e = v - v_k$.
\end{proposition}
The convergence rate is $O(\sqrt{\kappa})$; by Proposition~\ref{prop:shrinkage},
Ledoit-Wolf shrinkage compresses $\kappa$ monotonically in $\alpha$, directly reducing
iteration count. The requirement $\alpha > 0$ in Algorithm~\ref{alg:elim} ensures that each
active-submatrix $\hat\Sigma_{\mathrm{LW},\mathcal{A}}$ is strictly positive definite,
so Proposition~\ref{prop:cg} applies at every outer step and the inner CG solve is
guaranteed to converge to any prescribed tolerance in finitely many iterations.

On the S\&P~500 problem this brings CG iterations from 684 (sample
covariance, $\alpha=0$) down to 82 ($\alpha=0.5$, heavy-shrinkage range \cite{demiguel2009,ledoit2003jpm});
the oracle intensity $\alpha_{\mathrm{LW}} \approx 0.017$ reduces iterations
to 516 (one quarter), reflecting the Frobenius objective's insensitivity to
the near-zero noise eigenvalues that govern $\kappa$.

\paragraph{First-order baseline}
The benchmarks of Section~\ref{sec:results} include \emph{proximal gradient}: the
projected-gradient comparator of the companion paper~\cite[Section~5]{schmelzer2026nncg},
applied to the budget case with the $O(n\log n)$ simplex
projection~\cite{duchi2008,condat2016} and the same matrix-free gradient operator as CG.
It is loop-free and needs no dual feasibility check, but it is not a Krylov method and
converges at $O(\kappa)$ against CG's $O(\sqrt{\kappa})$ at the same $O(nT)$ per-step
cost, which makes it the slowest solver on ill-conditioned problems (the unshrunk
S\&P~500 system has $\kappa \approx 109{,}000$); only when shrinkage compresses $\kappa$
to $O(1)$ does its iteration count become competitive. The companion paper also
explains why the \emph{unaccelerated} method is the honest first-order comparator --
Nesterov momentum would match the $O(\sqrt\kappa)$ rate but not the Krylov subspace; a
FISTA variant~\cite{beck2009} empirically needs $1.5$--$1.8\times$ fewer
iterations than plain proximal gradient yet remains $7$--$10\times$ slower than CG on
the S\&P~500 benchmark at $\alpha=0.5$.

\subsection{Complexity Summary}\label{ssec:complexity}

Table~\ref{tab:complexity} collects the asymptotic cost of each method.
``Extra memory'' is storage beyond the $T \times n$ input matrix $X$; $\varepsilon$
is the solver tolerance; $s$ is the number of active-set outer steps.  The two
matrix-free methods never form $\hat\Sigma$ and require only $O(n)$ working memory
beyond $X$. For a general balance system the CG row scales by the $p+1$ right-hand
sides of Section~\ref{sec:pd}, which share one operator; the table shows the budget
case run in the experiments.

\begin{table}[ht]
\centering
\footnotesize
\begin{tabular}{lcccc}
\toprule
Method & Per-step cost & Iterations & Extra memory & Assembles $\hat\Sigma$ \\
\midrule
Clarabel (interior point)    & $O(n^3)^{\ddagger}$        & $O(\log 1/\varepsilon)$              & $O(n^2)$ & Yes \\
OSQP (operator splitting)    & $O(n^3)^{\ddagger}$ setup        & $O(1/\varepsilon)$                   & $O(n^2)$ & Yes \\
KKT (dense Cholesky) $+$ active set & $O(k^2 T)$ & $s^{\dagger}$ & $O(k^2)$ & Yes \\
CG matrix-free $+$ active set& $O(kT)$         & $s^{\dagger} \cdot O(\sqrt{\kappa})$ & $O(n)$   & No \\
Proximal gradient            & $O(nT)$         & $O(\kappa)$                          & $O(n)$   & No \\
\bottomrule
\end{tabular}
\caption{Asymptotic complexity per solver for $n$ assets, $T$ return observations,
and condition number $\kappa = \kappa(\hat\Sigma_{\mathrm{LW}})$.
$^\dagger$\,$s$ is the number of active-set outer steps (6--18 on S\&P~500 data;
see Table~\ref{tab:sp500}); the CG ``Iterations'' column counts $s \cdot O(\sqrt\kappa)$
total inner iterations, where each inner iteration is a pair of $O(kT)$ matrix-vector
products. Here $k \leq n$ is the active-set size; the outer loop starts at $k=n$
so the first inner solve dominates, giving effective $O(n^2)$ scaling empirically
(Section~\ref{sec:results}). The reported S\&P~500 CG counts (684, 82, 516) are
totals summed across all outer steps; because the first solve runs at the full
$k=n$ active set and dominates, they are governed by the full-system condition
number $\kappa(\hat\Sigma_{\mathrm{LW}})$ bounded in Proposition~\ref{prop:shrinkage}
and Corollary~\ref{cor:si}, not by the smaller $\kappa$ of the converged active
submatrix.
$^{\ddagger}$For the direct solvers the per-step column is the dominant cost of a
\emph{single factorisation}, not a cheap iteration: Clarabel refactorises an $O(n^3)$
dense KKT system at every interior-point step, whereas OSQP factorises once at
$O(n^3)$ setup and then reuses the cached factor in $O(n^2)$ back-solves per ADMM
step. The iteration-count entries are the standard worst-case rates
($O(\log 1/\varepsilon)$ for the interior-point method, the ergodic $O(1/\varepsilon)$
for ADMM); both are pessimistic in practice, where Clarabel converges in $11$--$16$
iterations and OSQP in roughly $75$--$225$ (Tables~\ref{tab:sp500}--\ref{tab:ftse}).
The table compares asymptotic structure, not the constant factors that determine
wall-clock time. The KKT row assembles the active-set Gram matrix and
Cholesky-factorises it ($O(k^2 T)$ assembly-dominated, $s$ outer steps, no inner
iterations); it is the dense counterpart to matrix-free CG and is faster when the
active set is small or the system well-conditioned (the crossover $\kappa < k^2$),
which is why it leads on the $86$-asset FTSE problem but trails CG as $n$ grows
(Section~\ref{sec:results}).}
\label{tab:complexity}
\end{table}
